\documentclass[10pt,a4paper]{amsart}
\usepackage[margin=19mm]{geometry}
\usepackage{amsmath,amssymb,mathtools}
\usepackage[scr=rsfs,cal=euler]{mathalfa}
\usepackage{xfrac}
\usepackage[unicode,hidelinks,bookmarks=false]{hyperref}
\newcommand{\ie}{i.e.\ }
\newcommand{\cf}{cf.\ }
\newcommand{\resp}{respectively\ }
\newcommand{\Subsection}{\subsection}
\providecommand{\nobreakdash}{\nobreak\hbox{-}}
\setlength{\emergencystretch}{2em}
\multlinegap0pt
\usepackage{xcolor}
\usepackage{algorithm}\usepackage{algpseudocode}
\usepackage{amssymb}
\usepackage{mathtools}
\usepackage[shortlabels]{enumitem}
\usepackage{tcolorbox}
\newtheorem{lemma}{Lemma}
\newtheorem{corollary}{Corollary}
\newcommand{\regret}{\mathfrak{R}}
\title{Bandit Convex Optimization with Gradient Prediction Adaptivity}
\author{Shuche Wang, Adarsh Barik, Vincent Y. F. Tan}
\date{Source: arXiv:2605.22191v1 (CC BY 4.0)}
\begin{document}
\maketitle

\noindent\emph{Source and licence.} Bandit Convex Optimization with Gradient Prediction Adaptivity. Version arXiv:2605.22191v1 (CC BY 4.0), distributed by arXiv under the Creative Commons Attribution 4.0 International licence, http://creativecommons.org/licenses/by/4.0/, as recorded in the arXiv licence metadata for this exact version and shown on the abstract page https://arxiv.org/abs/2605.22191v1. Full licence notice: \texttt{licenses/arxiv-2605.22191-CC-BY-4.0.txt}.\par\smallskip

\noindent\textit{Editorial scope.} Counted unit: the article's lemma lem:phase-residual-vs-Sbar (source0.tex lines 2105--2163). The article's complete original proof of it is delivered with it (source0.tex lines 2165--2406, one \texttt{proof} environment of 242 lines). No mathematics is written, repaired or re-derived: the statement and the proof are the article's own lines, in the article's own notation, and no retained line is substituted or re-flowed.  Retained uncounted as the source-local context the proof needs: lines 1428--1436; lines 1465--1477; lines 2073--2080.  Nothing was omitted from any retained span, so the package carries no omission record.  No retained line contains a citation, so the wrapper carries no bibliography and no bibliography entry.  No retained line contains a figure environment, an image-inclusion command or drawing code, so no image is delivered and the package declares no asset dependency.  Editorial formatting only, in addition: an AMS \texttt{amsart} preamble that restates the article's own declarations and macros used by the retained lines, namely the environments \texttt{lemma}, \texttt{corollary} on separate counters, together with the standard packages those lines need.  The retained environments print in this document as Lemma 1, Corollary 1 in order of appearance, whereas the article prints the same items with the numbering of its own full document.  The complete native source of the article as distributed in the arXiv e-print for this exact version is bundled unchanged as \texttt{sources/201118/source0.tex}.
\begin{lemma}[Second Moment Bound]
    \label{lem:second-moment-center} Assume $f_{t}$ is $\beta$-smooth. Then the
    second moment of the estimator error around the prediction $m_{t}$ is
    bounded by:
    \begin{equation*}
        \mathbb{E}\left[\|\hat g_{t}-m_{t}\|^{2}\mid \mathcal{H}_{t-1}\right]\le
        2d\|g_{t}^{y}-m_{t}\|^{2}+ \frac{d^{2}}{2}\beta^{2}\delta^{2}.
    \end{equation*}
\end{lemma}

\begin{corollary}[Second Moment Bound with Gradient at $x_{t}$]
    \label{cor:second-moment-play} Under the assumptions of Lemma~\ref{lem:second-moment-center},
    \begin{equation*}
        \mathbb{E}\left[\|\hat g_{t}-m_{t}\|^{2}\middle| \mathcal{H}_{t-1}\right
        ] \le 4d\,\mathbb{E}\left[\|g_{t}^{x}-m_{t}\|^{2}\;\middle|\; \mathcal{H}
        _{t-1}\right] + (d^{2}/2+4d)\beta^{2}\delta^{2}.
    \end{equation*}
    And in particular,
    \begin{equation*}
        \sum_{t=1}^{T}\mathbb{E}\left[\|\hat g_{t}-m_{t}\|^{2}\right]\le 4d\bar S
        _{T}+ (d^{2}/2+4d)\beta^{2}\delta^{2}T.
    \end{equation*}
\end{corollary}

\begin{lemma}
    \label{lem:finalS-phase} Let $R^{\mathrm{obs}}_{k}= \sum_{t\in I_k}\|\hat g
    _{t}-m_{t}\|^{2}$ denote the accumulated observable residual in phase $k$,
    and let $S_{k,E_k}$ be the final sensitivity budget used during phase $k$. Then, we have:
    \begin{equation*}
        S_{k,E_k}\le S_{\min} + \frac{R^{\mathrm{obs}}_{k}}{8d}.
    \end{equation*}
\end{lemma}

\begin{lemma}[Observed residual in a phase]
\label{lem:phase-residual-vs-Sbar}
Let $\regret_k^{\mathrm{obs}}=
\sum_{t\in I_k}\|\hat g_t-m_t\|^2$ be the accumulated observable residual in phase $k$. Suppose that the
sensitivity budgets in phase $k$ satisfy
\begin{equation*}
    S_{k,e}=S_{\min}2^{e-1},\qquad e=1,\ldots,E_k,
\end{equation*}
and that, within each epoch $(k,e)$, the perturbation radius is chosen as
\begin{equation*}
    \delta_{k,e}
    =
    \min\left\{
        \frac{\sqrt{S_{k,e}}}{d\beta H_k},
        \frac{1}{\beta\sqrt{H_k}},
        \frac r2
    \right\}.
\end{equation*}
Then, we have
\begin{equation*}
\mathbb E[\regret_k^{\mathrm{obs}}]
\le
\frac{
    4d\,\bar S_k+\frac{9}{2}\frac{S_{\min}}{H_k}
}{
    1-\frac{9}{16dH_k}
}.
\end{equation*}
In particular, since $d\ge 1$ and $H_k\ge 1$,
\begin{equation*}
\mathbb E[\regret_k^{\mathrm{obs}}]
\le
\frac{64}{7}d\,\bar S_k
+
\frac{72}{7}\frac{S_{\min}}{H_k}.
\end{equation*}
Consequently, the final sensitivity budget in phase $k$ satisfies
\begin{equation*}
\mathbb E[S_{k,E_k}]
\le
S_{\min}
+
\frac{1}{8d}
\cdot
\frac{
    4d\,\bar S_k+\frac{9}{2}\frac{S_{\min}}{H_k}
}{
    1-\frac{9}{16dH_k}
},
\end{equation*}
and, in particular,
\begin{equation*}
\mathbb E[S_{k,E_k}]
\le
\frac{8}{7}\bar S_k
+
\frac{16}{7}S_{\min}.
\end{equation*}
\end{lemma}

\begin{proof}
Let $Q_k=\mathbb E[\regret_k^{\mathrm{obs}}]$. From Corollary~\ref{cor:second-moment-play}, for every round $t\in I_k$, we have
\begin{equation*}
\mathbb E\!\left[
    \|\hat g_t-m_t\|^2
    \mid \mathcal H_{t-1}
\right]
\le
4d\|\nabla f_t(x_t)-m_t\|^2
+
(d^2/2+4d)\beta^2\delta_t^2,
\end{equation*}
where $\delta_t$ is the perturbation radius used at round $t$. Summing this
inequality over all rounds $t\in I_k$ and then taking total expectation gives
\begin{equation*}
\begin{aligned}
    Q_k
    &=
    \mathbb E\left[
        \sum_{t\in I_k}\|\hat g_t-m_t\|^2
    \right]  \\
    &\le
    4d\,
    \mathbb E\left[
        \sum_{t\in I_k}\|\nabla f_t(x_t)-m_t\|^2
    \right]
    +
    (d^2/2+4d)\beta^2
    \mathbb E\left[
        \sum_{t\in I_k}\delta_t^2
    \right]  \\
    &=
    4d\,\bar S_k
    +
    (d^2/2+4d)\beta^2
    \mathbb E\left[
        \sum_{t\in I_k}\delta_t^2
    \right].
\end{aligned}
\end{equation*}
We now decompose the last sum epoch by epoch. Since the perturbation radius is
fixed within each epoch $(k,e)$, we can write
\begin{equation*}
    \sum_{t\in I_k}\delta_t^2
    =
    \sum_{e=1}^{E_k}|J_{k,e}|\delta_{k,e}^2.
\end{equation*}
For every epoch $(k,e)$, by the definition of $\delta_{k,e}$,
\begin{equation*}
    \delta_{k,e}
    \le
    \frac{\sqrt{S_{k,e}}}{d\beta H_k}.
\end{equation*}
Therefore,
\begin{equation*}
\begin{aligned}
    (d^2/2+4d)\beta^2\delta_{k,e}^2
    &\le
    (d^2/2+4d)\beta^2
    \cdot
    \frac{S_{k,e}}{d^2\beta^2H_k^2}  \\
    &=
    \left(\frac12+\frac4d\right)
    \frac{S_{k,e}}{H_k^2}.
\end{aligned}
\end{equation*}
Since $d\ge1$, we have $\frac12+\frac4d\le \frac92$. Hence,
\begin{equation*}
    (d^2/2+4d)\beta^2\delta_{k,e}^2
    \le
    \frac92\frac{S_{k,e}}{H_k^2}.
\end{equation*}
Substituting this bound into the epoch-wise decomposition yields
\begin{equation*}
\begin{aligned}
    (d^2/2+4d)\beta^2
    \sum_{e=1}^{E_k}|J_{k,e}|\delta_{k,e}^2
    &\le
    \frac92
    \sum_{e=1}^{E_k}
    |J_{k,e}|\frac{S_{k,e}}{H_k^2}.
\end{aligned}
\end{equation*}
The budgets $S_{k,e}$ are nondecreasing in $e$, and hence
$S_{k,e}\le S_{k,E_k}$ for every $e\le E_k$. Moreover, the total number of
rounds in phase $k$ is at most $H_k$, so
\begin{equation*}
    \sum_{e=1}^{E_k}|J_{k,e}|=|I_k|\le H_k.
\end{equation*}
Therefore, pathwise,
\begin{equation*}
\begin{aligned}
    \sum_{e=1}^{E_k}
    |J_{k,e}|\frac{S_{k,e}}{H_k^2}
    &\le
    \sum_{e=1}^{E_k}
    |J_{k,e}|\frac{S_{k,E_k}}{H_k^2}  \\
    &=
    \frac{S_{k,E_k}}{H_k^2}
    \sum_{e=1}^{E_k}|J_{k,e}|  \\
    &\le
    \frac{S_{k,E_k}}{H_k}.
\end{aligned}
\end{equation*}
Combining the above estimates gives
\begin{equation*}
    Q_k
    \le
    4d\,\bar S_k
    +
    \frac{9}{2H_k}\mathbb E[S_{k,E_k}].
\end{equation*}

It remains to control the final sensitivity budget $S_{k,E_k}$. By
Lemma~\ref{lem:finalS-phase}, we have the pathwise bound
\begin{equation*}
    S_{k,E_k}
    \le
    S_{\min}
    +
    \frac{R_k^{\mathrm{obs}}}{8d}.
\end{equation*}
Taking expectation on both sides gives
\begin{equation*}
    \mathbb E[S_{k,E_k}]
    \le
    S_{\min}
    +
    \frac{Q_k}{8d}.
\end{equation*}
Substituting this inequality into the previous estimate for $Q_k$, we obtain
\begin{equation*}
\begin{aligned}
    Q_k
    \le
    4d\,\bar S_k
    +
    \frac{9}{2H_k}
    \left(
        S_{\min}
        +
        \frac{Q_k}{8d}
    \right)  =
    4d\,\bar S_k
    +
    \frac{9}{2}\frac{S_{\min}}{H_k}
    +
    \frac{9}{16dH_k}Q_k.
\end{aligned}
\end{equation*}
Moving the last term to the left-hand side gives
\begin{equation*}
    \left(
        1-\frac{9}{16dH_k}
    \right)Q_k
    \le
    4d\,\bar S_k
    +
    \frac{9}{2}\frac{S_{\min}}{H_k}.
\end{equation*}
Therefore,
\begin{equation*}
    Q_k
    \le
    \frac{
        4d\,\bar S_k+\frac{9}{2}\frac{S_{\min}}{H_k}
    }{
        1-\frac{9}{16dH_k}
    }.
\end{equation*}
This proves the first claimed bound.

Since $d\ge1$ and $H_k\ge1$, we have
\begin{equation*}
    1-\frac{9}{16dH_k}
    \ge
    1-\frac{9}{16}
    =
    \frac{7}{16}.
\end{equation*}
Thus,
\begin{equation*}
\begin{aligned}
    Q_k
    \le
    \frac{16}{7}
    \left(
        4d\,\bar S_k
        +
        \frac{9}{2}\frac{S_{\min}}{H_k}
    \right) =
    \frac{64}{7}d\,\bar S_k
    +
    \frac{72}{7}\frac{S_{\min}}{H_k}.
\end{aligned}
\end{equation*}
This proves the simplified residual bound.

Finally, applying Lemma~\ref{lem:finalS-phase} once more,
\begin{equation*}
    \mathbb E[S_{k,E_k}]
    \le
    S_{\min}
    +
    \frac{Q_k}{8d}.
\end{equation*}
Using the explicit bound on $Q_k$ gives
\begin{equation*}
\begin{aligned}
    \mathbb E[S_{k,E_k}]
    \le
    S_{\min}
    +
    \frac{1}{8d}
    \left(
        \frac{64}{7}d\,\bar S_k
        +
        \frac{72}{7}\frac{S_{\min}}{H_k}
    \right) =
    S_{\min}
    +
    \frac{8}{7}\bar S_k
    +
    \frac{9}{7dH_k}S_{\min}.
\end{aligned}
\end{equation*}
Since $dH_k\ge1$,
\begin{equation*}
    \frac{9}{7dH_k}S_{\min}
    \le
    \frac97S_{\min}.
\end{equation*}
Hence,
\begin{equation*}
    \mathbb E[S_{k,E_k}]
    \le
    \frac87\bar S_k
    +
    \frac{16}{7}S_{\min}.
\end{equation*}
This completes the proof.
\end{proof}

\end{document}
